\section{Conclusion}\label{sec:conclusion}
% Objective: Answer the question, state both contributions at their true size, name what would settle it.

No. An independent valuation built from public option quotes, a public rates
curve and the terms in the filing itself does not reproduce the issuer's
disclosed estimated value. It lands below it on all six notes, by \$11.41 to
\$27.92 per \$1{,}000 across the five whose barriers are covered by quoted
strikes, against simulation standard errors between \$0.24 and \$0.37.

The funding rate, which the filings themselves name as the input the reader
cannot see, does not account for it. Solving for the spread that would
reconcile the two valuations returns a negative number on every note, from
$-278$ to $-123$ basis points against an A-rated corporate benchmark of 69. The
sign is what carries the result: a disclosed value computed with a higher
internal funding rate would reconcile at a positive spread, and none of them
do.

Taking the remaining candidates one at a time excludes four more. Observation
timing moves values further away. The proxied Dow leg and the choice of term
extrapolation move the right way but are too small, or matter on too few notes,
to be the common cause. Suboptimal exercise accounts for part of three notes
and is arithmetically unavailable for the two five-year ones, whose disclosed
values sit above what the notes are worth with the issuer's call removed
altogether.

Two measured candidates survive on both direction and magnitude. The
correlation assumption would close each gap with a uniform increase of 0.15 to
0.27, which is large but runs the way a dealer's mark is known to run against a
realised estimate. The shape of the volatility surface would close it
comfortably: a fully flat surface overshoots the disclosed value on five of the
six notes. A third candidate cannot be measured with a local-volatility model
at all, and the published direction favours it: \citet{Farkas2022} report that
richer volatility dynamics price autocallables above local volatility, moving
the price while leaving the fit to European options untouched. None of the
three is published in the filing, and one of them, correlation, is not even
named in the lists of principal valuation inputs that two of these issuers
print.

\subsection*{What this paper contributes, and how much}

Two things, both narrower than the subject they sit in.

The first is a verification exercise on the product class the nearest published
work set aside. \citet{Bauer2020} reproduce issuer-disclosed values for German
retail certificates and leave autocallables to later work; this paper prices
six of them from public inputs and reports the funding-spread inversion as a
diagnostic, which converts an unexplained difference into a statement about
what the difference cannot be. The result is a case study on six notes from one
week, not an estimate of anything about the market. Whether these gaps are
typical is unknown, and nothing here should be read as suggesting they are.

The second is a census of what the exercise can reach at all. Of 1{,}086
supplements filed on Form 424B2 in the sample week, 129 are final worst-of
notes with a contingent coupon and a disclosed estimated value, and 71 of those
can be priced from public option quotes. A reader checking this disclosure can
therefore reach 55 per cent of the qualifying market, and the inaccessible
remainder is not a random sample of it: 33 of the 42 notes excluded for
unlisted options reference a single US sector index that has no options market
of its own. That bound is a fact about the disclosure regime rather than about
any issuer, and it would hold for any outside party attempting the same check.

\subsection*{What would settle it}

None of the three surviving candidates requires access to an issuer's models to
test. Two of them, though, are not cleanly separable from each other, which
shapes the order the tests should be run in.

Correlation could be tested directly with dispersion trade data or with quotes
on multi-asset structures from which an implied correlation can be backed out.
If dealer correlation marks on these underlyings sit 0.15 to 0.27 above the
realised estimate over the relevant horizon, correlation is sufficient on its
own and the question becomes how much of the gap it takes. If they sit much
closer, it is not.

The surface dynamics could be tested by repricing the same six notes under a
stochastic-local volatility model calibrated to the same quotes. The
single-asset version of that test has been run: \citet{Farkas2022} find a
Heston local-stochastic volatility model prices autocallables above local
volatility, with the stochastic parameters moving the price while the European
fit is untouched. What is missing is the multi-asset case, where
\citet{DeCol2017} find the correlations between stochastic variances to be
central, and where the dynamics candidate and the correlation candidate stop
being separable. Running it on these six notes would show whether the residual
reported here survives a model that can see the dimension this one cannot.

The sample is the third thing to extend, and the cheapest. The 129 final
supplements identified for one week would support a cross-sectional analysis of
the disclosed margin against tenor, coupon, barrier distance and the number of
underlyings, and the screening code that produced them runs forward on any
later week. Whether the pattern reported here survives contact with a hundred
notes rather than six is the question this paper cannot answer and the one most
worth asking next.
